\chapter{Hardware y alimentación}
\label{cap:hardware}

Los tres capítulos anteriores describen qué hace cada placa. Este describe el
soporte físico que las sostiene. De dónde sale la energía y cómo se reparte,
sobre qué van montadas, cómo llegan hasta ellas los mandos que el usuario toca
y qué las encierra. El reparto de tareas entre placas está en el
capítulo~\ref{cap:diseno}, la lectura de los canales del panel en el
apartado~\ref{sec:panel} y el tratamiento de las notas en el
apartado~\ref{sec:midi}.

\section{Componentes y equipos}
\label{sec:materiales}

La Tabla~\ref{tab:materiales} recoge los componentes del instrumento. No hay
ninguna pieza fabricada a medida salvo las cuatro piezas de la carcasa, en
línea con el condicionante de partida del apartado~\ref{sec:requisitos}. El
banco de trabajo se reduce a una impresora de filamento con cama de
220$\times$220~mm, una fuente de laboratorio, un soldador y un
multímetro. Los importes se recogen en el anexo de presupuesto.

\begin{table}[htbp]
    \centering
    \caption{Componentes del instrumento.}
    \label{tab:materiales}
    \begin{tabular}{llc}
        \toprule
        \textbf{Bloque} & \textbf{Componente} & \textbf{Cant.} \\
        \midrule
        Módulos
          & Daisy Seed 1.2 & 1 \\
          & CYD ESP32-2432S028R, pantalla táctil de 2,8 pulgadas & 1 \\
          & ESP32-S3 N16R8 & 1 \\
          & Shield \gls{midi} con optoacoplador 6N138 & 1 \\
        \midrule
        Alimentación
          & Alimentador externo de 230~V a 5~V / 3~A & 1 \\
          & Conector de alimentación de 2,1~mm, de chasis & 1 \\
          & Interruptor de encendido ON-OFF-ON & 1 \\
          & Condensador electrolítico de 470~µF, reserva del raíl & 1 \\
          & Condensador electrolítico de 100~µF, filtro de audio & 2 \\
          & Condensador cerámico de 100~nF & 8 \\
          & Resistencia fusible de 3,3~$\Omega$ y 1~W & 1 \\
          & Diodo Schottky 1N5817 & 1 \\
        \midrule
        Panel
          & Potenciómetro lineal de 10~k$\Omega$ & 6 \\
          & Conmutador ON-OFF-ON, selector de filtro & 1 \\
          & Resistencia de 10~k$\Omega$, divisor del selector & 2 \\
          & Pantalla \gls{oled} de 0,96 pulgadas, SSD1306 & 1 \\
          & Jack de audio de 3,5~mm mono, cuerpo de plástico & 1 \\
        \midrule
        Soporte
          & Placa perforada de tiras, 100$\times$160~mm, baquelita & 1 \\
          & Tira de pines hembra de 2,54~mm, 20 contactos & 2 \\
          & Tira de pines hembra de 2,54~mm, 22 contactos & 2 \\
          & Conector JST de 1,25~mm, enlace con la \gls{cyd} & 2 \\
          & Cable de silicona de 24~AWG & --- \\
          & Filamento PLA y tornillería M3 con tuerca & --- \\
        \bottomrule
    \end{tabular}
\end{table}

El teclado que aporta las notas no figura en la tabla. Por la decisión del
apartado~\ref{sec:arquitectura} es un aparato aparte, con su propia
alimentación, y el instrumento funciona con cualquier controlador \gls{midi}.

\section{La alimentación}
\label{sec:alimentacion}
\hypertarget{RE2.2c}{}\label{RE2.2c}

Las cuatro cargas cuelgan en paralelo de un único raíl de 5~V, que es la
tensión de entrada de los tres módulos y del shield, y cada placa regula
internamente sus propios 3,3~V. La Figura~\ref{fig:alimentacion} recoge el
esquema completo.

\begin{figure}[htbp]
    \centering
    \includegraphics[width=\textwidth]{alimentacion.pdf}
    \caption{Esquema de la alimentación. El condensador de reserva \texttt{C4}
             queda aguas arriba del interruptor y la rama del procesador de
             audio es la única con elemento en serie.}
    \label{fig:alimentacion}
\end{figure}

El consumo del conjunto se estimó en unos 600~mA, de modo que un alimentador de 3~A
deja margen de sobra. Va fuera de la carcasa, así que dentro no hay tensión 
de red en ningún punto. Alimentar todo desde un solo raíl tiene además un tiene además 
un efecto que el UART, el SPI y el MIDI necesitan: las masas de las tres placas quedan 
unidas.

\subsection*{El interruptor y el condensador de reserva}

El raíl no se alimenta directamente desde el conector de entrada. El
electrolítico de 470~µF (\texttt{C4}) queda aguas arriba y el interruptor
aguas abajo, de manera que al cerrar, con el condensador ya cargado, las placas
ven un flanco rápido en lugar de la rampa lenta con la que sube el alimentador.
La razón es que durante esa rampa las placas pasan un tiempo largo en la zona
en la que su circuito de reinicio no sabe decidir, y arrancan de forma
intermitente. Con el flanco rápido y el condensador aportando el pico de
corriente que piden al arrancar el problema desaparece. Para que esto sea
posible el raíl de 5~V deja de ser una tira continua y se corta en un punto,
aguas arriba de todas las cargas, es el interruptor de encendido el que decide cuando
se cierra el circuito y alimenta las cargas. El raíl de masa, en cambio, no se 
interrumpe nunca. Cortarlo obligaría a las corrientes de retorno a buscar camino
por los cables de datos.

\subsection*{Los filtros}

El raíl lo comparten la pantalla y dos microcontroladores conmutando. El ruido que
generan puede acabar en la sección analógica del procesador de audio y salir por el 
jack. Los condensadores de la Figura~\ref{fig:alimentacion} lo atacan aprovechando que
la impedancia de un condensador es un circuito abierto en continua y tiende a un cortocircuito
conforme sube la frecuencia. Este comportamiento se describe con la ecuación \eqref{eq:impedancia}, 
donde $j$ es la unidad imaginaria, $\omega$ la frecuencia angular y $C$ la capacidad del condensador.

\begin{equation}
    Z_C = \frac{1}{j\omega C},
    \label{eq:impedancia}
\end{equation}

Puesto en paralelo con la alimentación, un condensador deriva a
masa la componente de alta frecuencia sin tocar la continua que alimenta las
placas. Por eso el bulk son dos en paralelo y no uno: el electrolítico de
470~µF aporta la capacidad que sostiene los transitorios lentos y el cerámico
de 100~nF es el que es efectivo en la parte alta del espectro.

La rama del procesador de audio lleva además un elemento en serie. La resistencia
\texttt{R1} de 3,3~$\Omega$ junto con los 100~µF de \texttt{C2} forman un paso bajo 
de primer orden cuya frecuencia de corte es la ecuación~\eqref{eq:corte}, que es la 
frecuencia a la que la señal se atenúa 3~dB.

\begin{equation}
    f_c = \frac{1}{2\pi R C} \approx 482~\text{Hz},
    \label{eq:corte}
\end{equation}

De este modo, el ruido de conmutación llega atenuado a la entrada de alimentación 
de la placa. El cerámico \texttt{C3} cubre la parte del espectro que el electrolítico
ya no ve, y va soldado entre las patillas de entrada y de masa del módulo, que son 
adyacentes. Ese elemento en serie no puede ir en el raíl común porque su caída la 
pagarían también las demás cargas. Los parámetros del optoacoplador del shield
están especificados con una alimentación de 4,5~V \cite{Vishay6N138}, que es la
más baja a la que el fabricante garantiza algo, así que sobre un raíl de 5~V el
margen es estrecho para poner nada en serie.
La resistencia elegida es del tipo fusible, es decir, está diseñada para
abrirse ante una sobrecarga. Conviene saberlo porque si se abre, la placa de audio 
deja de recibir alimentación y el instrumento queda mudo.

\subsection*{El diodo de la rama de la pantalla táctil}

La \gls{cyd} se reprograma por su propio conector \gls{usb}, así que en algún
momento va a estar alimentada por las dos fuentes a la vez. El diodo Schottky
\texttt{D1} en su rama resuelve el conflicto. La placa toma la más alta de las 
dos fuentes, así la caída del Schottky sería suficiente para que la \gls{usb} predomine.
Se elige Schottky por su caída directa baja, del orden de tres décimas de
voltio, que sobre un raíl de 5~V es asumible. Los otros dos módulos llevan ese
diodo integrado y no lo necesitan.

\subsection*{Las masas}

El fabricante del procesador de audio exige unir la masa analógica con la
digital en toda aplicación \cite{DaisySeed}, y esa unión se hace en un único
punto y cerca de la placa. Conviene precisar qué significa eso, porque unidas en
un punto las dos masas son el mismo nodo y a corriente continua están al mismo
potencial. Lo que las distingue no es el nodo sino el camino: el cobre que
separa cualquier punto de masa del punto de unión tiene resistencia, y por él
circula la corriente de retorno de lo que haya aguas arriba. Medido sobre la
placa, entre una y otra masa hay unos 1,7~mV.

De ahí salen las dos decisiones de trazado. El manguito del jack se lleva a la
patilla de masa analógica y no al raíl general, para que el retorno de la salida
de audio no comparta cobre con el de los dos microcontroladores. Los
potenciómetros, en cambio, se referencian a la masa digital, que quedaba mucho
más a mano en el trazado, porque 1,7~mV sobre un fondo de escala de 3,3~V son un
0,05\,\% y no mueven una lectura de mando. La patilla superior sigue colgando de
la alimentación analógica, así que la cancelación que se explicó en el
apartado~\ref{sec:panel} queda intacta.

\section{Diseño del circuito y trazado de la placa}
\label{sec:veroboard}

El montaje definitivo va sobre una placa perforada de tiras de cobre de
100$\times$160~mm, que a paso de 2,54~mm son 39 tiras de 63 agujeros cada una.
El circuito se dibujó en Fritzing sobre su vista de placa de tiras, que es a la
vez esquema y plano de montaje. En una placa de este tipo la posición de
cada componente forma parte del circuito. El resultado es la
Figura~\ref{fig:fritzing}. Sobre ese fichero se verificó el trazado
decodificándolo hueco a hueco.

\begin{figure}[htbp]
    \centering
    \includegraphics[width=0.86\textwidth]{fritzing_veroboard.pdf}
    \caption{Trazado de la placa. Los dos módulos van sobre zócalos, la
             alimentación ocupa el extremo libre de los raíles y todo lo que el
             usuario toca sube al panel por el mazo.}
    \label{fig:fritzing}
\end{figure}

\subsection*{Los cortes de pista}

Un módulo con dos filas de
patillas, colocado perpendicular a las tiras, deja enfrentadas la patilla $p$ y
la $41-p$ sobre la misma tira, y por tanto en cortocircuito. Colocarlo paralelo
no es la solución, porque entonces las veinte patillas de una fila caen sobre
la misma tira. Hay que cortar, por tanto, las tiras que pasan bajo las placas. 
Hay que llevar especial cuidado con cinco de esas uniones ya que son destructivas 
si se olvidan. La peor pone la alimentación analógica directamente contra la masa
analógica, otra mete los 5~V de la entrada de alimentación en una \gls{gpio} y una tercera 
deja el cursor de un potenciómetro sobre la salida de audio.

Tras hacer los agujeros se comprueban con doble medida. Una a lo largo de la tira
cruzando el agujero, que tiene que quedar en abierto, y otra contra las dos tiras
vecinas, por si la broca ha raspado el borde. De hecho, gracias a la segunda medida 
se cazó una rebaba de cobre que había creado un cortocircuito.

\subsection*{Reparto de la placa}

Para la disposición en la placa se tomaron unos acuerdos.
La tira 1 es el raíl de masa y la tira 2 el de 5~V. El reparto a lo largo de
los 63 agujeros se hizo por naturaleza de la señal. El enlace \gls{spi} corre a
10~MHz, por lo que es la señal más agresiva de la placa y sus flancos se
acoplan con facilidad a lo que tenga cerca. Por este motivo su recorrido tiene
que ser corto y no pasar cerca del jack ni de los cursores.
Mientras que la salida de audio y la masa analógica se van al borde opuesto,
donde además caen las patillas de los potenciómetros. Así se acaba creando una zona
digital, una mixta y una analógica, y la sección de alimentación en el extremo libre
de los raíles, que es donde la obliga a estar el corte del interruptor.

Conviene subrayar que esta placa no es «el circuito» sino un «concentrador».
Los dos módulos sobre zócalos, la alimentación con su filtro y una fila de 
conexiones. Nada de lo que el usuario toca o mira va soldado a ella. Todo eso 
vive en el panel y baja por cables a la placa, como en un pedal de guitarra.

El trazado se calcó sobre la cara de cobre en lugar de sobre la de baquelita, y
como la vista de la herramienta es la cara de componentes, la placa salió en
espejo respecto al diseño. Con los cortes ya taladrados se optó por adaptarse a la
placa y volver a trazar esta vez contando sobre la baquelita. Por este motivo, en la placa
se pueden visualizar cortes de pista que no son necesarios y otros que sí lo son. 
La Figura~\ref{fig:fritzing} es la versión final. El cambio, además, obligó a reasignar
tres cursores y el selector de filtro a otras patillas.

\section{El panel de control}
\label{sec:panel-fisico}

La Figura~\ref{fig:panel} es el alzado del panel, dibujado a escala a partir
del mismo fichero paramétrico del que salió la pieza impresa.

\begin{figure}[htbp]
    \centering
    \includegraphics[width=0.72\textwidth]{panel.pdf}
    \caption{Alzado del panel frontal. El trazo discontinuo es el contorno de
             las dos placas, que quedan por detrás; el continuo, lo que se ve
             desde fuera. Cotas en milímetros.}
    \label{fig:panel}
\end{figure}

Los seis potenciómetros son lineales de 10~k$\Omega$ y cuelgan de la
alimentación analógica del procesador de audio, que es la disposición que el
fabricante recoge como aplicación típica \cite{DaisySeed}. El valor hace que
los seis sumen unos 2~mA de consumo sobre el raíl analógico. Mientras
que un valor menor lo cargaría sin motivo, valores mayores harían la lectura más
sensible al ruido captado por el cable. Cada cursor lleva un cerámico de 100~nF
a la masa, soldado en el extremo de la placa para poder filtrar lo que 
ha recogido el hilo. Las dos resistencias de 10~k$\Omega$ del divisor del selector,
cuya función se explicó en el apartado~\ref{sec:panel}, van también en la placa, 
así que al panel sólo llega un hilo de cada potenciómetro y uno del selector. La 
alimentación 3V3 analógica y la masa van en daisy chain, desde el selector hasta el 
último potenciómetro, y de ahí a la placa. 

Para el grosor del panel se tomó en cuenta la rosca del potenciómetro, que mide
7,0~mm, de la que hay que descontar arandela y tuerca. De esta manera, el máximo
que podía medir el panel de grosor eran 4,1~mm pero se optó por 3,4~mm porque era
más que suficiente. Se decidió también añadir el hueco para la patilla antigiro de
cada potenciómetro, que no es redonda sino una pletina rectangular de
2,8 por 1,2~mm, y que se corta pasante como el propio agujero de la rosca. Es
lo que impide que el cuerpo del potenciómetro gire dentro del panel al mover el
mando, con la tuerca limitándose a apretarlo contra la plancha.

Una parte complicada fueron las ventanas para las pantallas montadas en el panel.
La ventana de la \gls{cyd} lleva el labio izquierdo cuatro milímetros más adentro 
que el derecho para tapar un borde negro que asomaba. El labio, que apoya sobre el 
cristal en todo el perímetro, fue una fuente de pulsación fantasma como se comentó 
en el apartado \ref{sec:tactil}. Además, la pantalla va girada media 
vuelta por un motivo físico, que su conector \gls{usb} quede en el lado opuesto a la 
\gls{oled}. Esto hizo que se tuviera que rehacer parte de su firmware y recalibrarla 
entera, está descrito en el apartado~\ref{sec:tactil}.
La \gls{oled} se resolvió en tres niveles. Uno para el asiento de la placa, otro para
el hueco del cristal y un labio para acotar la parte visible de la pantalla con un canal 
para que su cable plano no quede pinzado. 

Para las etiquetas, se decidió que fuesen grabadas con un milímetro de
profundidad y que su tamaño se ajustase solo según el número de caracteres, para que
las largas no invadan a sus vecinas.

Fuera del panel quedan el jack, los conectores \gls{midi}, la entrada de
alimentación y el interruptor de encendido, que se llevaron a las paredes de la caja.

\section{La carcasa}
\label{sec:carcasa}

La carcasa no es una pieza de forma arbitraria, es un conjunto de agujeros a cotas
exactas y cada cota depende de una pieza física que hay que medir. Para solucionar esto, se
modeló con una herramienta de diseño paramétrico por texto, en la que corregir
una cota es cambiar un número, en lugar de con un modelador orgánico, donde
habría que rehacer la geometría. El método de trabajo fue ir imprimiendo probetas 
que comprueban si las cotas son correctas poco a poco. La probeta debía laminarse con el
mismo perfil que la pieza definitiva, ya que el diámetro real de un agujero
impreso depende de factores de la impresión como el ancho de línea, el número de perímetros 
o de la velocidad de impresión. Por eso la probeta se implementó dentro del propio fichero 
de la pieza, seleccionable con un parámetro, de modo que comparte parámetros con ella. 
Las probetas ayudan a depurar el modelo y a calibrar la impresora reduciendo el
tiempo de prueba y error y el desperdicio de filamento. La Figura~\ref{fig:carcasa} muestra 
la carcasa montada y despiece.

El cuerpo es una cuña de 44~mm de alto por delante y 62 por detrás, lo que deja
la cara inclinada a algo menos de nueve grados. Tiene una bahía trasera plana que
lleva la planta a 186 por 186,6~mm, y está partido en tres piezas, como muestra
la Figura~\ref{fig:carcasa}. Cada pieza se imprime en la orientación en la
que su geometría queda vertical o apoyada para minimizar el uso de soportes y para
ser más preciso con los contornos críticos. Por ejemplo, al imprimir la trasera tumbada sobre su cara exterior,
los huecos de los conectores \gls{midi} se convierten en agujeros verticales pasantes.
Así desaparecen a la vez los voladizos y los puentes de impresión que habrían empeorado
la precisión de la pieza.

Todas las uniones son por tornillo contra tuerca metálica embutida y ninguna con
rosca en el plástico. El material del que está impreso, PLA, no aguanta bien las 
fuerzas de rosca y puede deshacerse con el tiempo. Además, para todos los agujeros se 
crearon parámetros de holgura editables. Estos valores de holgura vendrán dados por lo
comentado anteriormente, probetas y configuraciones de la impresora. De hecho, los
agujeros impresos salen entre 0,2 y 0,35~mm menores que la cota real puesta en el modelo, 
así que, por ejemplo, los pasos de tornillo de 3~mm van dibujados entre 3,4 y 3,8.

La Figura~\ref{fig:instrumento} muestra el instrumento terminado, ya montado y
en funcionamiento.

\begin{figure}[H]
    \centering
    \includegraphics[width=\textwidth]{carcasa.pdf}
    \caption{La carcasa, renderizada desde el modelo paramétrico. A la
             izquierda el conjunto montado; a la derecha el despiece, con el
             panel separado de la bandeja y las dos piezas de la trasera.}
    \label{fig:carcasa}
\end{figure}

\begin{figure}[H]
    \centering
    \includegraphics[height=7.5cm]{instrumento_panel.jpg}\hfill
    \includegraphics[height=7.5cm]{instrumento_uso.jpg}
    \caption{El instrumento terminado. A la izquierda, el panel de control con
             los seis mandos, el selector de tipo de filtro y los huecos de las
             dos pantallas. A la derecha, en funcionamiento: el dedo fija el
             punto del espacio latente, la pantalla táctil dibuja la forma de
             onda que se genera en ese punto y el display de servicio muestra
             el parámetro que se está editando.}
    \label{fig:instrumento}
\end{figure}

Con la caja montada, el instrumento queda descrito por completo, del conjunto
de datos al plástico. Queda pendiente la comprobación eléctrica sobre la placa
ya soldada, que debe cerrar la cuestión de la tensión de entrada del procesador
de audio, y queda comprobar que el conjunto hace lo que se dice que hace, que
es el asunto del capítulo siguiente.
